% GENERATED FILE. Edit canonical lessons, then run npm run generate:books.
% canonical-source-hash: fnv1a64:1459c5995130542a
% canonical-lessons: MW-C252-had, MW-C252-sas, MW-C252-haphno, MW-C252-chabno, MW-C252-pir

\chapter{Bones and Breath}
\label{ch:mw-bones-and-breath}
\hlchaptermodality{252}

\begin{chapteropening}
\textbf{By the end of this chapter:} \emph{I can say where it aches and when I am out of breath.}

Complete the last lesson of chapter 252: I can say where it aches and when I am out of breath.
\end{chapteropening}

\section[\texorpdfstring{\mw{हाड} (hāḍ)}{हाड (hāḍ)}]{\mw{हाड} (hāḍ) — a bone}

\label{lesson:MW-C252-had}



\begin{quote}
Before the new one: say the Marwadi for an elbow, then the Marwadi for the waist.
\end{quote}

\subsection*{You'll want to know: \mw{हाड}}
\textbf{\mw{हाड}} — \emph{hāḍ} — \textquotedblleft{}a bone\textquotedblright{}.

\begin{quote}
\mw{म्हारो} \mw{हाड} \mw{कोनी} \mw{टूट्यो।}
\end{quote}

\begin{quote}
\emph{mhāro hāḍ konī ṭūṭyo.}
\end{quote}

\begin{quote}
My bone did not break.
\end{quote}

A masculine noun.

\begin{cousinweb}
\textbf{\mw{हाड}} (\emph{hāḍ}) is inherited from Sanskrit \emph{haḍḍa}, \textquotedblleft{}bone\textquotedblright{}; CDIAL lists Old Marwari \emph{hāḍī} under it (CDIAL 13952).
\end{cousinweb}

\begin{grammarlens}[title={What you've built}]
One more word for this chapter.
\end{grammarlens}

\subsection*{Your turn}
\begin{itemize}
\raggedright
  \item \emph{Say it:} \emph{hāḍ}
  \item \emph{Say it:} \emph{hāḍ}, once more
  \item \emph{Say it:} \emph{mhāro hāḍ konī ṭūṭyo.} — My bone did not break.
  \item \emph{Recall:} say the Marwadi for an elbow, then the Marwadi for the waist, then say \emph{hāḍ} again
\end{itemize}

\begin{tcolorbox}[breakable,colback=teal!4,colframe=teal!35!black,title={Before you move on}]
What does \mw{हाड} mean? (\textquotedblleft{}A bone\textquotedblright{}.) Say it once more.
\end{tcolorbox}
\section[\texorpdfstring{\mw{सांस} (sā̃s)}{सांस (sā̃s)}]{\mw{सांस} (sā̃s) — breath}

\label{lesson:MW-C252-sas}



\begin{quote}
Before the new one: say the Marwadi for the waist, then the Marwadi for a bone.
\end{quote}

\subsection*{You'll want to know: \mw{सांस}}
\textbf{\mw{सांस}} — \emph{sā̃s} — \textquotedblleft{}breath\textquotedblright{}.

\begin{quote}
\mw{लांबी} \mw{सांस} \mw{लो।}
\end{quote}

\begin{quote}
\emph{lā̃bī sā̃s lo.}
\end{quote}

\begin{quote}
Take a long breath.
\end{quote}

A feminine noun.

\begin{cousinweb}
\textbf{\mw{सांस}} (\emph{sā̃s}) is inherited from Sanskrit \emph{śvāsa}, \textquotedblleft{}breath\textquotedblright{}; CDIAL lists Marwari \emph{sā̃so}, \textquotedblleft{}panting\textquotedblright{}, and Hindi \emph{sā̃s}, \textquotedblleft{}breath\textquotedblright{}, under it (CDIAL 12769).
\end{cousinweb}

\begin{grammarlens}[title={What you've built}]
One more word for this chapter.
\end{grammarlens}

\subsection*{Your turn}
\begin{itemize}
\raggedright
  \item \emph{Say it:} \emph{sā̃s}
  \item \emph{Say it:} \emph{sā̃s}, once more
  \item \emph{Say it:} \emph{lā̃bī sā̃s lo.} — Take a long breath.
  \item \emph{Recall:} say the Marwadi for the waist, then the Marwadi for a bone, then say \emph{sā̃s} again
\end{itemize}

\begin{tcolorbox}[breakable,colback=teal!4,colframe=teal!35!black,title={Before you move on}]
What does \mw{सांस} mean? (\textquotedblleft{}Breath\textquotedblright{}.) Say it once more.
\end{tcolorbox}
\section[\texorpdfstring{\mw{हांफणो} (hā̃phṇo)}{हांफणो (hā̃phṇo)}]{\mw{हांफणो} (hā̃phṇo) — to pant; to be out of breath}

\label{lesson:MW-C252-haphno}



\begin{quote}
Before the new one: say the Marwadi for a bone, then the Marwadi for breath.
\end{quote}

\subsection*{You'll want to know: \mw{हांफणो}}
\textbf{\mw{हांफणो}} — \emph{hā̃phṇo} — \textquotedblleft{}to pant; to be out of breath\textquotedblright{}.

\begin{quote}
\mw{डूंगर} \mw{माथै} \mw{दादा} \mw{हांफै} \mw{है।}
\end{quote}

\begin{quote}
\emph{ḍū̃gar māthai dādā hā̃phai hai.}
\end{quote}

\begin{quote}
On the hill, grandfather is out of breath.
\end{quote}

Past (he …): \mw{वो} \mw{छोरो} \mw{हांफ्यो}, \emph{vo chhoro hā̃phyo}.

\begin{cousinweb}
\textbf{\mw{हांफणो}} (\emph{hā̃phṇo}) is inherited from a reconstructed \emph{*hamph}, \textquotedblleft{}pant\textquotedblright{}. CDIAL gives no Marwari form in this entry; it lists Hindi \emph{hā̃phnā}, \textquotedblleft{}to pant\textquotedblright{} (CDIAL 13973).
\end{cousinweb}

\begin{grammarlens}[title={What you've built}]
One more word for this chapter.
\end{grammarlens}

\subsection*{Your turn}
\begin{itemize}
\raggedright
  \item \emph{Say it:} \emph{hā̃phṇo}
  \item \emph{Say it:} \emph{hā̃phṇo}, once more
  \item \emph{Say it:} \emph{ḍū̃gar māthai dādā hā̃phai hai.} — On the hill, grandfather is out of breath.
  \item \emph{Recall:} say the Marwadi for a bone, then the Marwadi for breath, then say \emph{hā̃phṇo} again
\end{itemize}

\begin{tcolorbox}[breakable,colback=teal!4,colframe=teal!35!black,title={Before you move on}]
What does \mw{हांफणो} mean? (\textquotedblleft{}To pant; to be out of breath\textquotedblright{}.) Say it once more.
\end{tcolorbox}
\section[\texorpdfstring{\mw{चाबणो} (chābṇo)}{चाबणो (chābṇo)}]{\mw{चाबणो} (chābṇo) — to chew}

\label{lesson:MW-C252-chabno}



\begin{quote}
Before the new one: say the Marwadi for breath, then the Marwadi for to pant.
\end{quote}

\subsection*{You'll want to know: \mw{चाबणो}}
\textbf{\mw{चाबणो}} — \emph{chābṇo} — \textquotedblleft{}to chew\textquotedblright{}.

\begin{quote}
\mw{रोटी} \mw{चाबो।}
\end{quote}

\begin{quote}
\emph{roṭī chābo.}
\end{quote}

\begin{quote}
Chew your flatbread.
\end{quote}

Past (I …): \mw{म्हैं} \mw{चणा} \mw{चाब्या}, \emph{mhaĩ chaṇā chābyā}.

\begin{cousinweb}
\textbf{\mw{चाबणो}} (\emph{chābṇo}) is inherited from Sanskrit \emph{carvati}, \textquotedblleft{}chews\textquotedblright{}; CDIAL lists Marwari \emph{cābṇo} under it (CDIAL 4711).
\end{cousinweb}

\begin{grammarlens}[title={What you've built}]
One more word for this chapter.
\end{grammarlens}

\subsection*{Your turn}
\begin{itemize}
\raggedright
  \item \emph{Say it:} \emph{chābṇo}
  \item \emph{Say it:} \emph{chābṇo}, once more
  \item \emph{Say it:} \emph{roṭī chābo.} — Chew your flatbread.
  \item \emph{Recall:} say the Marwadi for breath, then the Marwadi for to pant, then say \emph{chābṇo} again
\end{itemize}

\begin{tcolorbox}[breakable,colback=teal!4,colframe=teal!35!black,title={Before you move on}]
What does \mw{चाबणो} mean? (\textquotedblleft{}To chew\textquotedblright{}.) Say it once more.
\end{tcolorbox}
\section[\texorpdfstring{\mw{पीड़} (pīṛ)}{पीड़ (pīṛ)}]{\mw{पीड़} (pīṛ) — an ache; a deep pain}

\label{lesson:MW-C252-pir}



\begin{quote}
Before the new one: say the Marwadi for to pant, then the Marwadi for to chew.
\end{quote}

\subsection*{You'll want to know: \mw{पीड़}}
\textbf{\mw{पीड़}} — \emph{pīṛ} — \textquotedblleft{}an ache; a deep pain\textquotedblright{}.

\begin{quote}
\mw{म्हारी} \mw{पीठ} \mw{मांय} \mw{पीड़} \mw{है।}
\end{quote}

\begin{quote}
\emph{mhārī pīṭh mā̃y pīṛ hai.}
\end{quote}

\begin{quote}
I have an ache in my back.
\end{quote}

A feminine noun.

\begin{cousinweb}
\textbf{\mw{पीड़}} (\emph{pīṛ}) is inherited from Sanskrit \emph{pīḍā}, \textquotedblleft{}pain\textquotedblright{}; CDIAL lists Old Marwari \emph{pīḍa} under it (CDIAL 8227).
\end{cousinweb}

\begin{grammarlens}[title={What you've built}]
One more word for this chapter.
\end{grammarlens}

\subsection*{Your turn}
\begin{itemize}
\raggedright
  \item \emph{Say it:} \emph{pīṛ}
  \item \emph{Say it:} \emph{pīṛ}, once more
  \item \emph{Say it:} \emph{mhārī pīṭh mā̃y pīṛ hai.} — I have an ache in my back.
  \item \emph{Recall:} say the Marwadi for to pant, then the Marwadi for to chew, then say \emph{pīṛ} again
\end{itemize}

\begin{tcolorbox}[breakable,colback=teal!4,colframe=teal!35!black,title={Before you move on}]
What does \mw{पीड़} mean? (\textquotedblleft{}An ache; a deep pain\textquotedblright{}.) Say it once more.
\end{tcolorbox}
